% How accurate LIGHT v0_5 is: v0_4 with the 5s -> 5p electron excitation
% on Jung et al.'s measured cross sections, and a Kr-buffered e-beam
% laser held out. No number here is typed by hand. They come from
%   overview_v0_5_fast.tex, _slow.tex, _unc.tex   make_overview_v0_5.py  (\wg)
%   system_map_v0_5_numbers.tex                   make_system_map_v0_5.py (\wm)
%   overview_v0_4_fast.tex, _slow.tex, _unc.tex   make_overview_v0_4.py  (\wf)
%   overview_v0_3_fast.tex, _slow.tex             make_overview_v0_3.py  (\wk)
%   overview_v0_2_fast.tex, _slow.tex             make_overview_v0_2.py  (\wi)
%   energy_ladder_v0_5.tex                        make_energy_ladder.py  (\lad)
% Build from LIGHT/docs:
%   python3 make_system_map_v0_5.py && python3 make_overview_v0_5.py
%   python3 make_energy_ladder.py v0_5
%   latexmk -pdf accuracy_overview_v0_5.tex
\documentclass[11pt]{article}
\usepackage[a4paper,margin=2.2cm]{geometry}
\usepackage[T1]{fontenc}
\usepackage[scaled=0.95]{helvet}
\renewcommand{\familydefault}{\sfdefault}
\usepackage{amsmath}
\usepackage{sansmath}
\usepackage{booktabs,tabularx,array}
\usepackage{xcolor}
\usepackage{float}
\usepackage{enumitem}
\usepackage{needspace}
\usepackage{tikz}
\usepackage{graphicx}
\usetikzlibrary{arrows.meta,decorations.pathmorphing}
\setlist{nosep, leftmargin=1.2em}
\setlength{\parindent}{0pt}
\setlength{\parskip}{5pt}
\definecolor{ink2}{HTML}{52514E}
\renewcommand{\topfraction}{0.95}
\renewcommand{\textfraction}{0.03}

\input{overview_v0_2_fast}
\input{overview_v0_2_slow}
\input{overview_v0_3_fast}
\input{overview_v0_3_slow}
\input{overview_v0_4_fast}
\input{overview_v0_4_slow}
\input{overview_v0_4_unc}
\input{system_map_v0_5_numbers}
\input{overview_v0_5_fast}
\input{overview_v0_5_slow}
\input{overview_v0_5_unc}
\input{energy_ladder_v0_5}
\makeatletter
\newcommand\wi[1]{\ifcsname wi@#1\endcsname\csname wi@#1\endcsname
  \else\errmessage{no v0_2 number '#1'. Rerun make_overview_v0_2.py}\fi}
\newcommand\wk[1]{\ifcsname wk@#1\endcsname\csname wk@#1\endcsname
  \else\errmessage{no v0_3 number '#1'. Rerun make_overview_v0_3.py}\fi}
\newcommand\wf[1]{\ifcsname wf@#1\endcsname\csname wf@#1\endcsname
  \else\errmessage{no v0_4 number '#1'. Rerun make_overview_v0_4.py}\fi}
\newcommand\wm[1]{\ifcsname wm@#1\endcsname\csname wm@#1\endcsname
  \else\errmessage{no map number '#1'. Rerun make_system_map_v0_5.py}\fi}
\newcommand\wg[1]{\ifcsname wg@#1\endcsname\csname wg@#1\endcsname
  \else\errmessage{no v0_5 number '#1'. Rerun make_overview_v0_5.py}\fi}
\newcommand\lad[1]{\ifcsname lad@#1\endcsname\csname lad@#1\endcsname
  \else\errmessage{no ladder number '#1'. Rerun make_energy_ladder.py}\fi}
\makeatother

% chemistry and units
\newcommand{\unit}[1]{\ensuremath{\,\mathrm{#1}}}
\newcommand{\el}{\ensuremath{\mathrm{e}^-}}
\newcommand{\Krs}{Kr\ensuremath{^*}}
\newcommand{\Krpp}{Kr\ensuremath{^{**}}}
\newcommand{\Krp}{Kr\ensuremath{^+}}
\newcommand{\Krtp}{Kr\ensuremath{_2^+}}
\newcommand{\Krts}{Kr\ensuremath{_2^*}}
\newcommand{\KtF}{Kr\ensuremath{_2}F\ensuremath{^*}}
\newcommand{\Ft}{F\ensuremath{_2}}
\newcommand{\Fm}{F\ensuremath{^-}}
\newcommand{\ra}{\ensuremath{\to}\ }
\newcommand{\Edep}{\ensuremath{E_\mathrm{dep}}}
\newcommand{\gIP}{\ensuremath{g_0I_\mathrm{sat}/P_\mathrm{in}}}
\newcommand{\x}[1]{\ensuremath{\times}#1}
\newcommand{\cm}[1]{\ensuremath{\,\mathrm{cm^{#1}s^{-1}}}}
\newcommand{\pc}[1]{\wm{#1}\,\%}
\newcommand{\Lam}{\ensuremath{\text{{\unsansmath$\Lambda$}}}}
\newcommand{\pctof}[2][0]{\pgfmathparse{100*#2}%
  \pgfmathprintnumber[fixed,precision=#1]{\pgfmathresult}}
% a model-over-measurement factor as a signed percent error
\newcommand{\errof}[1]{\pgfmathparse{100*(#1-1)}%
  \pgfmathprintnumber[fixed,precision=0,showpos]{\pgfmathresult}}
% model value over measured value as a signed percent error
\newcommand{\errratio}[2]{\pgfmathparse{100*(#1/#2-1)}%
  \pgfmathprintnumber[fixed,precision=0,showpos]{\pgfmathresult}}

\begin{document}
\sansmath

{\LARGE\bfseries How accurate LIGHT v0\_5 is}\par
\vspace{4pt}
{\large The basic network of LIGHT v0\_5, version \lad{version}, on 11
October 2026: v0\_4 with the 5s $\to$ 5p electron excitation on the
cross sections Jung et al.\ measured, and a krypton-buffered e-beam laser
held out, with its energy ladder and its checks}\par
\vspace{10pt}

\noindent\fbox{\parbox{\dimexpr\linewidth-2\fboxsep-2\fboxrule\relax}{%
\setlength{\parskip}{5pt}%
\textbf{Bottom line.} LIGHT v0\_5 reproduces the fitted KrF
decays~\cite{eden78,gadomski92} within \wg{sc.train.hi}\,\% rms
(\wg{sc.train.lo} to \wg{sc.train.hi}\,\%) and the decays no fit
used~\cite{xu93} within \wg{sc.xu_figs5_7}\,\% rms, as v0\_2 to v0\_4 do.
At the Tahoma laser~\cite{tahoma89} it gives the small-signal gain to
\errof{\wg{tah.kin.g0.x}}\,\% (\wg{tah.kin.g0} against the measured
\wg{tah.g0.meas}\unit{\%/cm}; v0\_4 \errof{\wf{tah.kin.g0.x}}\,\%), the
absorption to \errof{\wg{tah.kin.alpha.x}}\,\%, the net gain to
\errof{\wg{tah.kin.net.x}}\,\%, the saturation intensity to
\errof{\wg{tah.isat.x}}\,\% and the formation efficiency to
\errof{\wg{tah.eta.x}}\,\%. No fit used a laser's numbers and no rate was
set by one.

What v0\_5 changes is one rate: the electron pumping of the 5p rung,
which v0\_4's overview named as deciding the partition, runs on the cross
sections Jung et al.\ measured out of the metastable levels of
Kr~\cite{jung05,jung06} in place of a 1976 fit~\cite{kannari85,jacob76}.
The measurement confirms the fit: its Maxwellian rate is
\x{\wg{jung.ratio.116}} the fit's at 1.16\unit{eV} (Tahoma's electrons),
\x{\wg{jung.ratio.3}} at 3\unit{eV} and \x{\wg{jung.ratio.58}} at
5.8\unit{eV}, and the 248-nm share at the reference point moves from
\wm{om.v04} to \wm{out.248}\,\% of \Edep{} (\wf{ref.248} to
\wg{ref.248}\,\% under the kinetic EEDF).

Its own uncertainty, propagated from its sources (Section~\ref{sec:accuracy}:
\wg{unc.n} inputs at the two ends their sources allow, in quadrature), is
$\pm$\wg{unc.ref.quad.pct}\,\% on the 248-nm share and
$\pm$\wg{unc.g0.quad.pct}, $\pm$\wg{unc.alpha.quad.pct} and
$\pm$\wg{unc.net.quad.pct}\,\% on Tahoma's gain, absorption and net gain,
the largest term everywhere but the absorption \wg{unc.ref.top1}, a band
with no measurement behind it. A second laser now bears on that term:
Kleikamp et al.'s krypton-buffered KrF laser~\cite{kleikamp86} peaks at
\wg{kl.meas.kr}\unit{atm} of Kr, and LIGHT, scored with the paper's own
saturated balance, puts the optimum at \wg{kl.kr.sat.mid}\unit{atm}
(\errratio{\wg{kl.kr.sat.mid}}{\wg{kl.meas.kr}}\,\%; \wg{kl.kr.sat.lo} to
\wg{kl.kr.sat.hi} over the pump's uncertainty), because at 2\unit{atm} it
quenches the threshold inversion \x{\wg{kl.k.ratio}} as fast as the
laser's balance allows. Nothing is retuned.

The largest errors: Hakuta et al.'s C-A decays~\cite{hakuta90},
\wg{sc.hakuta_fig2} and \wg{sc.hakuta_fig3}\,\% rms; Burnham and Searles'
KrF(B) lifetime~\cite{burnham77}, \wg{sc.bs_tau_b}\,\%; the VUV yield of
pure krypton at low pump~\cite{morozov08}, \errof{\wk{pure.vuv.x}}\,\%
(v0\_3's run); Tahoma's absorption, \errof{\wg{tah.kin.alpha.x}}\,\%,
\wg{tah.k2f.share}\,\% of it \KtF{} at the cross section Hakuta et al.\
measured~\cite{hakuta87}; and the Kr pressure optimum of Kleikamp et
al.'s laser, \errratio{\wg{kl.kr.sat.mid}}{\wg{kl.meas.kr}}\,\%, with its
F$_2$ optimum (\wg{kl.meas.f2}\unit{mbar}) not reproduced at all.}}
\medskip

LIGHT v0\_5 is v0\_4~\cite{lightv04} with one row changed and two data
sets added. K76, \el{} + Kr(5s) \ra \el{} + Kr(5p), was Kannari's Table I
fit $A\,T_e^{1/2}\exp(-C/T_e)$ to e-beam laser data; it is now the lumped
cross section built from the ten measured excitations out of each
metastable level (Section~\ref{sec:changes}), under both closures, with
the fit kept as a switch. The TALIF review of Gazeli et al.~\cite{gazeli21rev}
supplies a measured 5p$'$ lifetime that cross-checks the rung's
radiative rate. Kleikamp et al.'s laser~\cite{kleikamp86} is a second
device, held out. Everything else, the 5p rung, the optics, KrF(B,C) + F,
the ion-ion limbs, the walls, the wave and the kinetic EEDF, is v0\_4's,
and v0\_4's overview describes it. The \wg{sc.same} kinetic datasets have
no electrons, so v0\_5 scores them exactly as v0\_2, and
\texttt{make\_overview\_v0\_5.py} refuses to write a number that differs.
Table~\ref{tab:score} lists every test.

\begin{table}[!tbp]
\centering\footnotesize
\renewcommand{\arraystretch}{1.12}
\begin{tabularx}{\linewidth}{@{}>{\raggedright\arraybackslash}X
  >{\raggedright\arraybackslash}p{2.3cm}>{\raggedright\arraybackslash}p{2.4cm}
  >{\raggedright\arraybackslash}p{2.3cm}>{\raggedright\arraybackslash}p{1.9cm}@{}}
\toprule
test & measured & v0\_5 & v0\_4 & role of the data\\
\midrule
energy ledger closure & & below \wg{wl.closure} with walls,
  \wg{wv.closure} with the wave & the same & bookkeeping\\
KrF(B) decay against Kr and F$_2$, Eden 1978 & & rms \wg{sc.eden_kr} and
  \wg{sc.eden_f2}\,\% & the same & train\\
KrF(B,C) slow decay, Gadomski 1992 Figs.\ 1 and 2 & & rms
  \wg{sc.ksu_fig1} and \wg{sc.ksu_fig2}\,\% & the same & train\\
C/B emission ratio at 300\unit{K}, Yu and Setser 1990 & \wi{ys.meas}
  $\pm$ \wi{ys.err} & \wg{sc.ys_ir300}\,\% & the same & train\\
KrF(B,C) decays, Xu 1993 Figs.\ 5 to 7 & & rms \wg{sc.xu_figs5_7}\,\% &
  the same & held out\\
decay constants Xu 1993 printed, 50 and 100 Torr Kr & \wi{xu.50.meas}
  and \wi{xu.100.meas}\unit{s^{-1}} & \x{\wg{xu.50}} and \x{\wg{xu.100}} &
  the same & structure check\\
C-A decays, Hakuta 1990 & & rms \wg{sc.hakuta_fig2} and
  \wg{sc.hakuta_fig3}\,\% & the same & held out, in conflict\\
$\tau_B$, Burnham and Searles 1977 & \wi{sc.bs.meas} $\pm$
  \wi{sc.bs.err}\unit{ns} & \wg{sc.bs_tau_b}\,\% & the same & held out, in
  conflict\\
C/B ratio 230\unit{K} over 300\unit{K}, Yu and Setser & \wi{ysr.meas}
  $\pm$ \wi{ysr.err} & \wg{sc.ys_ratio_230_300}\,\% & the same & held out\\
\midrule
Tahoma small-signal gain $g_0$, kinetic EEDF & \wg{tah.g0.meas}\unit{\%/cm} &
  \wg{tah.kin.g0}, \x{\wg{tah.kin.g0.x}} & \wf{tah.kin.g0}, \x{\wf{tah.kin.g0.x}} &
  device, score only\\
\quad the same, v0\_2's closure & & \wg{tah.default.g0},
  \x{\wg{tah.default.g0.x}} & \wg{tah.v04.g0}, \x{\wg{tah.v04.g0.x}} &
  device, score only\\
Tahoma absorption $\alpha_0$, kinetic EEDF & \wg{tah.alpha.meas}\unit{\%/cm}
  & \wg{tah.kin.alpha}, \x{\wg{tah.kin.alpha.x}} & \wf{tah.kin.alpha},
  \x{\wf{tah.kin.alpha.x}} & device, score only\\
Tahoma net gain $g_0-\alpha_0$, kinetic EEDF & \wg{tah.net.meas}\unit{\%/cm}
  & \wg{tah.kin.net}, \x{\wg{tah.kin.net.x}} & \wf{tah.kin.net},
  \x{\wf{tah.kin.net.x}} & device, score only\\
Tahoma loss in pure Kr, 257.3\unit{nm} & \wg{tah.pure.meas}\unit{\%/cm} &
  \wg{tah.pure.alpha}, \x{\wg{tah.pure.x}} (248-nm cross sections) & the
  same & rough check\\
Tahoma saturation intensity $I_\mathrm{sat}$, v0\_2's closure &
  \wg{tah.isat.meas} $\pm$ \wi{tah.isat.err}\unit{MW/cm^2} & \wg{tah.isat},
  \x{\wg{tah.isat.x}} & \wf{tah.isat}, \x{\wf{tah.isat.x}} & device, score
  only\\
Tahoma formation efficiency \gIP, v0\_2's closure & \wg{tah.eta.meas} &
  \wg{tah.eta}, \x{\wg{tah.eta.x}} & \wf{tah.eta}, \x{\wf{tah.eta.x}} &
  device, score only\\
\midrule
Kr pressure optimum of Kleikamp 1986's Kr/F$_2$ laser, F$_2$
  \wg{kl.meas.f2}\unit{mbar} & \wg{kl.meas.kr}\unit{atm} &
  \wg{kl.kr.sat.mid}\unit{atm} (\wg{kl.kr.sat.lo} to \wg{kl.kr.sat.hi} over
  the pump's spread), the paper's saturated balance & not scored &
  device, held out\\
F$_2$ pressure optimum, Kr 2\unit{atm} & \wg{kl.meas.f2}\unit{mbar} &
  above \wg{kl.f2.top}\unit{mbar} (still rising) & & device, held out\\
three-body KrF$^*$ + Kr quenching their balance deduces, cm$^6$/s &
  $\wg{kl.meas.kR}$ (model-dependent) & $\wg{kl.k.trimer}$ at the
  window's \wg{kl.tg}\unit{K} & & device, score only\\
\bottomrule
\end{tabularx}
\caption{LIGHT v0\_5 against the data, beside v0\_4, every band at its
default unless a row says otherwise. A percentage is the deviation from
the data, signed for a single value and rms for a set. A \x{} factor is
the model over the measurement. ``The same'' is checked: the generator
stops if v0\_5's number differs from v0\_2's. The Kleikamp rows are new
(Section~\ref{sec:kleikamp}).}
\label{tab:score}
\end{table}

\section{Overall accuracy and its uncertainty}
\label{sec:accuracy}

\textbf{Against the data.} Table~\ref{tab:score} scores every test;
Table~\ref{tab:errors} reads the same numbers as errors, the model's
deviation from the measurement in percent of the measurement (rms over a
set of points), and names the largest.

\begin{table}[H]
\centering\footnotesize
\renewcommand{\arraystretch}{1.1}
\begin{tabularx}{\linewidth}{@{}>{\raggedright\arraybackslash}X
  >{\raggedright\arraybackslash}p{3.0cm}>{\raggedright\arraybackslash}p{3.0cm}
  >{\raggedright\arraybackslash}p{3.4cm}@{}}
\toprule
quantity & measured & v0\_5 & error\\
\midrule
KrF(B,C) decays the fit used, Eden 1978 and Gadomski 1992~\cite{eden78,gadomski92}
  & four sets & & within \wg{sc.train.hi}\,\% rms (\wg{sc.train.lo} to
  \wg{sc.train.hi})\\
KrF(B,C) decays no fit used, Xu 1993~\cite{xu93} & three figures & & within
  \wg{sc.xu_figs5_7}\,\% rms\\
C/B emission ratio at 300\unit{K}~\cite{yu90} & \wi{ys.meas} $\pm$ \wi{ys.err} & &
  \wg{sc.ys_ir300}\,\%\\
its 230\unit{K}/300\unit{K} ratio~\cite{yu90} & \wi{ysr.meas} $\pm$ \wi{ysr.err} & &
  \wg{sc.ys_ratio_230_300}\,\%\\
decay constants, 50 and 100\unit{Torr} Kr~\cite{xu93} & \wi{xu.50.meas},
  \wi{xu.100.meas}\unit{s^{-1}} & \x{\wg{xu.50}}, \x{\wg{xu.100}} &
  \errof{\wg{xu.50}}, \errof{\wg{xu.100}}\,\%\\
Tahoma small-signal gain $g_0$~\cite{tahoma89} & \wg{tah.g0.meas}\unit{\%/cm} &
  \wg{tah.kin.g0} & \errof{\wg{tah.kin.g0.x}}\,\%\\
Tahoma absorption $\alpha_0$ & \wg{tah.alpha.meas}\unit{\%/cm} & \wg{tah.kin.alpha} &
  \errof{\wg{tah.kin.alpha.x}}\,\%\\
Tahoma net gain $g_0-\alpha_0$ & \wg{tah.net.meas}\unit{\%/cm} & \wg{tah.kin.net} &
  \errof{\wg{tah.kin.net.x}}\,\%\\
Tahoma saturation intensity & \wg{tah.isat.meas} $\pm$ \wi{tah.isat.err}\unit{MW/cm^2}
  & \wg{tah.isat} & \errof{\wg{tah.isat.x}}\,\%\\
Tahoma formation efficiency \gIP{} & \wg{tah.eta.meas} & \wg{tah.eta} &
  \errof{\wg{tah.eta.x}}\,\%\\
Kleikamp 1986's Kr pressure optimum~\cite{kleikamp86} & \wg{kl.meas.kr}\unit{atm}
  & \wg{kl.kr.sat.mid}\unit{atm} & \errratio{\wg{kl.kr.sat.mid}}{\wg{kl.meas.kr}}\,\%\\
\midrule
\multicolumn{4}{@{}l}{\emph{the largest errors}}\\
C-A decays, Hakuta 1990~\cite{hakuta90} (in conflict with the fitted sets) & two
  figures & & \wg{sc.hakuta_fig2} and \wg{sc.hakuta_fig3}\,\% rms\\
KrF(B) lifetime, Burnham and Searles 1977~\cite{burnham77} (in conflict) &
  \wi{sc.bs.meas} $\pm$ \wi{sc.bs.err}\unit{ns} & & \wg{sc.bs_tau_b}\,\%\\
VUV yield of pure Kr at low pump, Morozov 2008~\cite{morozov08}, v0\_3's run &
  \wk{pure.vuv.meas} $\pm$ \wk{pure.vuv.err} & \wk{pure.vuv} &
  \errof{\wk{pure.vuv.x}}\,\%\\
Kleikamp 1986's Kr pressure optimum & \wg{kl.meas.kr}\unit{atm} &
  \wg{kl.kr.sat.mid}\unit{atm} & \errratio{\wg{kl.kr.sat.mid}}{\wg{kl.meas.kr}}\,\%\\
Kleikamp 1986's F$_2$ pressure optimum & \wg{kl.meas.f2}\unit{mbar} & above
  \wg{kl.f2.top} & not reproduced\\
Tahoma absorption $\alpha_0$ & \wg{tah.alpha.meas}\unit{\%/cm} & \wg{tah.kin.alpha}
  & \errof{\wg{tah.kin.alpha.x}}\,\%\\
\bottomrule
\end{tabularx}
\caption{The scores of Table~\ref{tab:score} as errors. A single value's
error is signed, model less measurement over the measurement; a set's is
the rms over its points. Tahoma's are under the kinetic EEDF, the
Kleikamp optimum under v0\_2's closure; the rest have no electrons.}
\label{tab:errors}
\end{table}

\textbf{Propagated from the sources.} As in v0\_4's overview, each input
is set to the two ends its source allows and the outputs recomputed
(Table~\ref{tab:unc}, under v0\_2's closure): the \wg{unc.nband} bands to
their ranges (Table~\ref{tab:bands}; k\_5s5p is now a scale on the
measured cross section, \wg{k.5s5p.lo} to \wg{k.5s5p.hi}, where in v0\_4
it ran from the fit to 0); the \wg{unc.nrow} constants the model uses at a
measured value whose source quotes an error and that no band covers,
\Krtp{} + \el{} ($\pm$10\,\%~\cite{shiu77}), \Krp{} + 2Kr
($\pm$9\,\%~\cite{grossl81}) and the \KtF{} lifetime ($\pm$9\,\%~\cite{xu93});
and the \wg{unc.nsig} cross sections with a quoted error analytically,
$\sigma_\mathrm{se}$ $\pm$\wg{unc.sig.se}\,\%~\cite{tellinghuisen76,kannari85}
and $\sigma$(\KtF) $\pm$\wg{unc.sig.k2f}\,\%~\cite{hakuta87}. Half the
difference between an input's two ends is its term; the quadrature sum
treats the terms as independent and each as one standard deviation, which
a band's range is not, so it is indicative; the linear sums of the
excursions are the ceiling on the spread, not a bound. Not propagated: the
electron-impact cross sections, now including Jung et al.'s, whose
$\pm$\wg{jung.unc.1s5} to \wg{jung.unc.1s3}\,\% the band k\_5s5p carries
instead, and the fitted KrF(B,C) set, which the training decays pin at
\wg{sc.train.hi}\,\% rms.

On the 248-nm share at the reference point (\wg{unc.ref.y0}\,\% of \Edep)
the quadrature sum is $\pm$\wg{unc.ref.quad} points
($\pm$\wg{unc.ref.quad.pct}\,\%; v0\_4 $\pm$\wf{unc.ref.quad.pct}\,\%), the
excursions sum to \wg{unc.ref.down} and \wg{unc.ref.up} points, and the
largest terms are \wg{unc.ref.top}. On Tahoma's gain
(\wg{unc.g0.y0}\unit{\%/cm}) it is $\pm$\wg{unc.g0.quad.pct}\,\%, the
largest terms \wg{unc.g0.top}; the measurement is \wg{unc.g0.dev.pct}\,\%
away, \wg{unc.g0.dev.quad} quadrature units. On the absorption
(\wg{unc.alpha.y0}\unit{\%/cm}) it is $\pm$\wg{unc.alpha.quad.pct}\,\%, the
largest terms \wg{unc.alpha.top}; the measurement is
\wg{unc.alpha.dev.pct}\,\% away, \wg{unc.alpha.dev.quad} units. On the net
gain (\wg{unc.net.y0}\unit{\%/cm}) it is $\pm$\wg{unc.net.quad.pct}\,\%, the
measurement \wg{unc.net.dev.pct}\,\% away, \wg{unc.net.dev.quad} units.
Every Tahoma measurement lies within one quadrature unit, and the
measured cross section has taken the 5s $\leftrightarrow$ 5p term from
v0\_4's $\pm$\wf{unc.ref.quad}-point sum down to its share here; what
remains is the trimer's temperature law, on which Section~\ref{sec:kleikamp}
now has a laser to say something.

\begin{table}[H]
\centering\scriptsize
\resizebox{\linewidth}{!}{%
\begin{tabular}{@{}llllll@{}}
\toprule
input & its two ends & 248\unit{nm} (\% of \Edep) & $g_0$ (\%/cm) & $\alpha_0$ (\%/cm) & $g_0-\alpha_0$\\
\midrule
\wgUncTable
\bottomrule
\end{tabular}}
\caption{The propagated uncertainty's inputs, each alone at its two ends,
the others at their defaults, and the four outputs at those ends (low end
to high end), under v0\_2's closure. \wg{unc.min} minutes of runs.}
\label{tab:unc}
\end{table}

\section{What v0\_5 changes}
\label{sec:changes}

\textbf{The measured cross section.} Jung, Stone, Boffard, Anderson and
Lin~\cite{jung06} measured the electron-impact excitation cross sections
out of the two metastable levels of Kr, 1s$_5$ ($J=2$) and 1s$_3$ ($J=0$),
into the \wg{jung.n.levels} levels of the 4p$^5$5p configuration, at 2 to
10\unit{eV} for every level and to 300\unit{eV} for the two largest
(their Tables II to IV; the Letter~\cite{jung05} prints the 3.5-eV
column). The totals at 6\unit{eV} are $\wg{jung.sum.1s5}\unit{cm^2}$ out
of 1s$_5$ and $\wg{jung.sum.1s3}$ out of 1s$_3$, with stated
uncertainties of $\pm$\wg{jung.unc.1s5} and $\pm$\wg{jung.unc.1s3}\,\%.
LIGHT's \Krs{} is the whole 4p$^5$5s manifold (the Biagi group at
9.915\unit{eV}) and its \Krpp{} the whole 5p manifold, so one lumped
cross section is needed: the statistical-weight mean over the four 5s
levels of each level's sum over the ten 5p levels. Nothing is measured
out of the two $J=1$ resonance levels, which at 1\unit{atm} are
radiation-trapped and populated with the metastables; they are given the
set of the metastable with the same ion core, their thresholds shifted to
their own energies. That is the one assumption, and the band k\_5s5p
brackets it: \wg{k.5s5p.lo} is the metastables alone (the $J=1$ levels
contributing nothing, half the manifold's weight) and \wg{k.5s5p.hi} the
measurement's upper error. The lumped cross section is
$\wg{jung.sig6}\unit{cm^2}$ at 6\unit{eV} (the metastables alone
$\wg{jung.sig6.meta}$, \wg{jung.meta.share} of it), linear between the
tabulated energies, from 0 at each pair's own threshold, and above the
tables on the shape of the largest column, 1s$_5\to$ 2p$_9$. The electron
loses the ledger's gap, \wm{gap5p}\unit{eV}, as in v0\_4, so the ledger
closes. Under v0\_2's closure the row is the Maxwellian table of that
cross section and its reverse the Klein-Rosseland table; under the
kinetic EEDF the population process carries the cross section itself.
Checks: the solver's forward rate for a Maxwellian matches the table to
\wg{pop.k} and the reverse to \wg{pop.sup}, and detailed balance over the
gap holds to \wg{pop.db} (Section~\ref{sec:checks}).

\textbf{The finding.} The measured set confirms the fit it replaces,
Kannari's Table I [76], Jacob and Mangano's $A\,T_e^{1/2}\exp(-C/T_e)$
with $A=\wg{k.76.A}$, $C=\wg{k.76.C}\unit{eV}$: the lumped rate is
\x{\wg{jung.ratio.1}} the fit's at 1\unit{eV}, \x{\wg{jung.ratio.116}} at
1.16, \x{\wg{jung.ratio.2}} at 2, \x{\wg{jung.ratio.3}} at 3 and
\x{\wg{jung.ratio.58}} at 5.8\unit{eV}, above it where the laser's
electrons sit and below it at the reference point's. So the partition
v0\_4 found rests on a measurement now and moves little:
Table~\ref{tab:ref} and Table~\ref{tab:tahoma}. \texttt{Params.k76 =
"kannari"} keeps the fit as a switch and reproduces v0\_4 to four digits.

\textbf{A TALIF lifetime.} The TALIF review of Gazeli et
al.~\cite{gazeli21rev} gives the natural lifetime of the Kr 5p$'$[3/2]$_2$
calibration level as \wg{tau.talif}\unit{ns}; PRAD's statistical mean of
the ten measured level lifetimes is \wg{tau5p}\unit{ns}, and Dzier\.z\k{e}ga
et al.~\cite{dzierzega00} give that level \wg{tau.2p2}\unit{ns}. No
change; the review tabulates no Kr self-quenching, so KPS stays on Gazeli
et al.'s Physics of Plasmas measurement~\cite{gazeli21}.

\textbf{A second laser.} Section~\ref{sec:kleikamp}.

\textbf{What v0\_4 added}, and this version keeps unchanged: the 5p rung
(KPS at $\wg{k.KPS}\cm{3}$, PRAD at \wg{tau5p}\unit{ns}, K27 at
$\wg{k.K27}\cm{3}$, the stepwise ionizations as bands at 0), KrF(B,C) + F
as a band, the 248-nm optics on Kannari's absorbers with Hakuta et al.'s
measured \KtF{} cross section, ion-ion recombination on Flannery's
turnover with the Langevin limb of the measured mobilities, the walls of
a gap and the EM wave with Turnbull et al.'s inverse bremsstrahlung, the
5s $\leftrightarrow$ 5p collisions, the ions' Coulomb collisions and the
wave's field term in the Boltzmann equation. v0\_4's overview~\cite{lightv04}
describes and checks them; their numbers here are unchanged from it
(the wave's \wg{wv.em}\unit{J\,cm^{-3}} beside the \wg{wv.given} given, the
walls' \wg{wl.share}\,\%, the field term's atom part to \wg{hf.atoms} and
ion part to \wg{hf.ions}, \texttt{ib.py} held by \wg{ib.checks} checks).

\section{How the numbers were made}

Every number in this note is written by \texttt{docs/make\_overview\_v0\_5.py}
from v0\_5 and from \texttt{ML/}, or by \texttt{docs/make\_system\_map\_v0\_5.py}
(the reference point under v0\_2's closure, its bands and its optics,
which the system map draws), and a test fails when those files no longer
match the model. v0\_4's numbers come from its own generated files,
unchanged. ML runs v0\_5 when \texttt{LIGHT\_VERSION} names it. Tahoma
runs under the prototype beam drive, a test fixture in \texttt{docs/} and
not part of the model, in ML's configuration (v0\_2's: the default network
with every band at its default and $n_e$ free from a seed of
\wi{tah.seed}\unit{cm^{-3}}). Kleikamp et al.'s laser runs under the same
drive at their pump through \texttt{docs/krf\_data/kleikamp\_v0\_5.py}. The
kinetic EEDF takes minutes a run, so Tahoma's variants, its saturation
intensity, its formation efficiency and the Kleikamp scans run on v0\_2's
closure (the kinetic closure moves v0\_5's Tahoma numbers by about 1\,\%,
Table~\ref{tab:score}); the kinetic numbers are the reference point and
Tahoma's default, \wg{ref.min} and \wg{tah.kin.min} minutes each.

\section{Tahoma and the optics}

\begin{table}[H]
\centering\footnotesize
\setlength{\tabcolsep}{3.5pt}
\resizebox{\linewidth}{!}{%
\begin{tabular}{@{}llllll@{}}
\toprule
v0\_2's closure & $g_0$ (\%/cm) & $\alpha_0$ (\%/cm) & $g_0-\alpha_0$ & $T_e$ (eV) & $n_e$ (cm$^{-3}$)\\
\midrule
measured~\cite{tahoma89} & \wg{tah.g0.meas} & \wg{tah.alpha.meas} & \wg{tah.net.meas} & & \\
\wgTahVariants
v0\_4 & \wg{tah.v04.g0} (\x{\wg{tah.v04.g0.x}}) & \wg{tah.v04.alpha}
  (\x{\wg{tah.v04.alpha.x}}) & \wg{tah.v04.net} (\x{\wg{tah.v04.net.x}}) &
  \wg{tah.v04.te} & \wg{tah.v04.ne}\\
v0\_3 & \wg{tah.v03.g0} (\x{\wg{tah.v03.g0.x}}) & \wg{tah.v03.alpha}
  (\x{\wg{tah.v03.alpha.x}}) & \wg{tah.v03.net} (\x{\wg{tah.v03.net.x}}) &
  \wg{tah.v03.te} & \wg{tah.v03.ne}\\
\midrule
v0\_5, kinetic EEDF & \wg{tah.kin.g0} (\x{\wg{tah.kin.g0.x}}) &
  \wg{tah.kin.alpha} (\x{\wg{tah.kin.alpha.x}}) & \wg{tah.kin.net}
  (\x{\wg{tah.kin.net.x}}) & \wg{tah.kin.te} & \wg{tah.kin.ne}\\
v0\_4, kinetic EEDF & \wf{tah.kin.g0} (\x{\wf{tah.kin.g0.x}}) &
  \wf{tah.kin.alpha} (\x{\wf{tah.kin.alpha.x}}) & \wf{tah.kin.net}
  (\x{\wf{tah.kin.net.x}}) & \wf{tah.kin.te} & \wf{tah.kin.ne}\\
\bottomrule
\end{tabular}}
\caption{Tahoma over the \wi{tah.win.lo} to \wi{tah.win.hi}\unit{ns}
window: v0\_5 with every band at its default and each variant alone, v0\_4
and v0\_3 (v0\_4's cross sections on their densities), and the kinetic
EEDF. The absorption by species with the defaults: \wg{tah.by}\unit{\%/cm}
(v0\_2's closure), \wg{tah.kin.by} (kinetic); pure Kr, no \Ft:
\wg{tah.pure.by}.}
\label{tab:tahoma}
\end{table}

At Tahoma's \wg{tah.default.te}-eV electrons the measured cycle is
\x{\wg{jung.ratio.116}} the fit's, cools the electrons a little further
(v0\_4's \wg{tah.v04.te}\unit{eV}) and lifts the gain by
\errratio{\wg{tah.default.g0}}{\wg{tah.v04.g0}}\,\%; the absorption does
not move, \wg{tah.k2f.share}\,\% of it being \KtF{}. The band's two ends,
the metastables alone and the measurement's upper error, move the gain
between \x{\wg{tah.k5s5plo.g0.x}} and \x{\wg{tah.k5s5phi.g0.x}}; the cycle
off gives \x{\wg{tah.nocycle.g0.x}}. The bands that move the gain most are
v0\_4's: the trimer's temperature law (\x{\wg{tah.tri.g0.x}}), the ion-ion
yield (\x{\wg{tah.iiyield.g0.x}}), KrF(B,C) + F at Kannari's analogue
(\x{\wg{tah.krff.g0.x}}), the ion-ion C share (\x{\wg{tah.iifc.g0.x}}) and
\Krts{} + F (\x{\wg{tah.kr2f.g0.x}}), none with a measurement behind the
band. In the window under the kinetic EEDF LIGHT's mean energy is the
kinetic one at \wg{tah.kin.reached} update, held by a field of at most
\wg{tah.kin.en}\unit{Td}, and the two energy balances agree to
\wg{tah.kin.gap} of the beam's electron power. ML's formation run gives
$g_0$ \wg{tah.ml.g0}\unit{\%/cm} with its own channels, $I_\mathrm{sat}$
\wg{tah.isat}\unit{MW/cm^2} and \gIP{} $=$ \wg{tah.eta}.

\section{The reference point}

At the notes' reference point, \wm{p}\unit{atm}, \pc{xF2} \Ft{},
\wm{P}\unit{kW\,cm^{-3}} for \wm{pulse}\unit{\mu s} and run to
\wm{tend}\unit{\mu s}, v0\_5 gives these shares of the deposited energy:

\begin{table}[H]
\centering\footnotesize
\begin{tabular}{@{}lrrrr@{}}
\toprule
 & v0\_5, kinetic & v0\_5, v0\_2's closure & v0\_4, kinetic & v0\_4, v0\_2's closure\\
\midrule
248\unit{nm}, KrF B-X & \wg{ref.248} & \wm{out.248} & \wf{ref.248} & \wg{v04.248}\\
272\unit{nm}, KrF B-A & \wg{ref.272} & \wm{out.272} & \wf{ref.272} & \wg{v04.272}\\
275\unit{nm}, KrF C-A & \wg{ref.275} & \wm{out.275} & \wf{ref.275} & \wg{v04.275}\\
420\unit{nm}, \KtF{} & \wg{ref.420} & \wm{out.420} & \wf{ref.420} & \wg{v04.420}\\
147\unit{nm}, \Krts{} & \wg{ref.147} & \wm{out.147} & \wf{ref.147} & \wg{v04.147}\\
near IR, 5p \ra 5s & \wg{ref.IR} & \wm{out.IR} & \wf{ref.IR} & \wg{v04.IR}\\
all bands & \wg{ref.all} & \wm{notes.emit} & \wf{ref.all} & \wg{v04.all}\\
heat, reactions & \wg{ref.heat} & \wm{out.heat} & \wf{ref.heat} & \wg{v04.heat}\\
to \Krs{}, gross & \wg{ref.krstar} & \wm{notes.krs} & \wf{ref.krstar} & \wg{v04.krstar}\\
to \Krp{} & \wg{ref.krplus} & \wm{notes.krp} & \wf{ref.krplus} & \wg{v04.krplus}\\
KrF(B) formed & \wg{ref.krfb} & \wm{krf.eff} & \wf{ref.krfb} & \wg{v04.krfb}\\
elastic heating & \wg{ref.elastic} & \wm{out.elastic} & \wf{ref.elastic} & \wg{v04.elastic}\\
$g_0$, 500 to 1000\unit{ns} (\%/cm) & \wg{ref.g0} & \wm{op.g0} & \wf{ref.g0} & \wg{v04.g0}\\
$\alpha$ (\%/cm) & \wg{ref.alpha} & \wm{op.alpha} & \wf{ref.alpha} & \wg{v04.alpha}\\
$T_g$ at the end (K) & \wg{ref.tg} & \wm{TgEnd} & \wf{ref.tg} & \wg{v04.tg}\\
$T_e$, pulse mean (eV) & \wg{ref.te} & \wm{Te} & \wf{ref.te} & \wg{v04.te}\\
\bottomrule
\end{tabular}
\caption{The reference point, percent of \Edep{} unless a row says
otherwise. v0\_5's v0\_2-closure column is the system map's run; v0\_4's
is the same run with Kannari's fit in K76's place.}
\label{tab:ref}
\end{table}

What moves the result is still the 5s $\leftrightarrow$ 5p cycle, on the
measurement now: the electrons pump \pc{brs.76} of the \Krs{} formed up to
the rung (\wg{ref.brs76}\,\% under the kinetic EEDF; v0\_4
\wf{ref.brs76}\,\%), the rung relaxes \pc{brp.KPS} of it on Kr and the
\wm{gap5p}\unit{eV} of every cycle heats the gas, \pc{cyc.e} of the input.
Under v0\_2's closure the electrons sit at \wm{Te}\unit{eV}, where the
measured rate is \x{\wm{k.76.ratio}} the fit's, so the share falls from
v0\_4's \wm{om.v04} to \wm{out.248}\,\% and the \Krts{} route gains it;
under the kinetic EEDF they sit at \wg{ref.te}\unit{eV}, where it is
\x{\wg{jung.ratio.3}}, and the share rises from \wf{ref.248} to
\wg{ref.248}\,\%. With the rung neutralised the share is \wm{om.v03}\,\%,
v0\_3's. Under the kinetic EEDF the field carries \wg{ref.ratio.lo} to
\wg{ref.ratio.hi} of the input through the pulse's second half and the two
energy balances agree to \wg{ref.gap}; the run takes \wg{ref.solutions}
kinetic solutions, \wg{ref.on} in the pulse, with $E/N$ between
\wg{ref.en.lo} and \wg{ref.en.hi}\unit{Td}, and closes its ledger below
\wg{ref.closure}.

Every band's two ends under v0\_2's closure, one band at a time
(\texttt{run\_v0\_5.py --bands}):

\begin{table}[H]
\centering\scriptsize
\resizebox{\linewidth}{!}{%
\begin{tabular}{@{}lllllllr@{}}
\toprule
band & default & range & 248\unit{nm} & 147\unit{nm} & $g_0$ & $\alpha$ & $T_g$ (K)\\
\midrule
\wmBands
\bottomrule
\end{tabular}}
\caption{Each band alone at its two ends, the others at their defaults
(\wm{bd.n} bands, \wm{bd.nzero} with no measurement and a default of 0).
The columns are the 248-nm and 147-nm shares of \Edep{} (\%), the
500-to-1000-ns gain and absorption (\%/cm) and the final gas
temperature, low end to high end. With every band at its default the
248-nm share is \wm{out.248}\,\%.}
\label{tab:bands}
\end{table}

\section{A second laser, held out: Kleikamp et al.'s Kr/F$_2$ laser}
\label{sec:kleikamp}

Kleikamp, Peters and Witteman~\cite{kleikamp86} pumped a KrF laser with
krypton as the only buffer in a coaxial e-beam diode (280\unit{kV},
\wg{kl.i}\unit{A/cm^2} on average, \wg{kl.pulse}-ns pulses, electrons of
at most \wg{kl.emax}\unit{keV} in the gas, a 1-cm tube 20\unit{cm} long,
an outcoupling of $\wg{kl.gamma}\unit{cm^{-1}}$). With F$_2$ at
\wg{kl.meas.f2}\unit{mbar} the output peaks at about \wg{kl.meas.kr}\unit{atm}
of Kr (their Fig.\ 3); with Kr at 2\unit{atm} at about
\wg{kl.meas.f2}\unit{mbar} of F$_2$ (Fig.\ 4); and their steady-state
cavity balance deduces a three-body KrF$^*$ + Kr quenching constant of
$\wg{kl.meas.kR}\cm{6}$ (their Table I: ``the quenching constants of Kr
and Ne are not known in the literature''). It is a held-out device, a
score and never a target.

LIGHT runs under the beam drive at their conditions. The deposited power
is not printed in the paper: from the collision stopping power of Kr at
170\unit{keV} in Berger and Seltzer's tables~\cite{berger64}, the source
the paper names, read to $\pm$15\,\%, it is 3.8\unit{MW\,cm^{-3}} at
2\unit{atm}; from their own pumping rate $B_\mathrm{Kr}$ with the W-value
of Kr, 5.6. The scans run at the mean, \wg{kl.P}\unit{MW\,cm^{-3}} at
2\unit{atm} proportional to the density, and again at \wg{kl.P.lo} and
\wg{kl.P.hi} times it. LIGHT has no cavity, so two proxies for the output
are compared with the measured optimum's location over the pulse's second
half (\wg{kl.win.lo} to \wg{kl.win.hi}\unit{ns}): the net small-signal gain
$g_0-\alpha_0$, and the paper's own balance of a saturated laser (its
Eqs.\ (1) to (3)): with the inversion clamped at threshold,
$[\mathrm{KrF}]_\mathrm{thr}=(\gamma+\alpha)/\sigma_\mathrm{s}$, the output
per length is $\gamma/(\gamma+\alpha)\,(P-Q\,[\mathrm{KrF}]_\mathrm{thr})$
with $P$ LIGHT's KrF(B) formation rate, $Q$ its loss rate per molecule
(quenching and radiation) and $\alpha$ its absorption. The optimum of the
second balances the pumping, which rises with the density, against the
quenching of the threshold inversion, which rises faster (the trimer goes
as the density squared), so it tests the trimer's rate at the laser's gas
temperature, the largest term of the propagated uncertainty.

\begin{table}[H]
\centering\footnotesize
\begin{tabular}{@{}lllllll@{}}
\toprule
Kr (atm) & $g_0$ (\%/cm) & $\alpha_0$ & $g_0-\alpha_0$ & $Q$ (s$^{-1}$) & output, saturated balance & $T_g$ (K)\\
\midrule
\wgKleikamp
\bottomrule
\end{tabular}
\caption{LIGHT v0\_5 at Kleikamp et al.'s conditions, F$_2$ at
\wg{kl.meas.f2}\unit{mbar}, the nominal pump, v0\_2's closure, the
\wg{kl.win.lo} to \wg{kl.win.hi}-ns window. The saturated-balance output
is relative to its maximum. The measured output peaks at about
\wg{kl.meas.kr}\unit{atm}.}
\label{tab:kleikamp}
\end{table}

The saturated balance peaks at \wg{kl.kr.sat.mid}\unit{atm}
(\wg{kl.kr.sat.lo} at \wg{kl.P.lo} times the pump, \wg{kl.kr.sat.hi} at
\wg{kl.P.hi}), the small-signal net gain at \wg{kl.kr.ss.mid}; the laser
peaked at \wg{kl.meas.kr}. At 2\unit{atm} LIGHT's KrF(B) loss rate is
$\wg{kl.Q}\unit{s^{-1}}$ at the window's $T_g$ of \wg{kl.tg}\unit{K} ($T_e$
\wg{kl.te}\unit{eV}, $n_e$ \wg{kl.ne}\unit{cm^{-3}}), the trimer most of it
at $\wg{kl.k.trimer}\cm{6}$ (G5's 300-K rate on the $(T_g/300)^{-3.32}$
law), \x{\wg{kl.k.ratio}} the constant the paper's balance deduces. With
G5 and G5C scaled by \wg{kl.trimer.scale}, which brings the effective
quenching to $\wg{kl.scaled.k}$, the saturated optimum moves to
\wg{kl.scaled.sat}\unit{atm} (the small-signal one to \wg{kl.scaled.ss});
that is a sensitivity run and not adopted, the laser's constant being a
model-dependent number (their $Q$ is three-body only, their output
insensitive to the outcoupling) that bounds the term rather than fixes
it. The F$_2$ optimum is not reproduced: both proxies still rise at
\wg{kl.f2.top}\unit{mbar} where the laser peaked at \wg{kl.meas.f2},
because LIGHT's absorption falls with F$_2$ (\KtF{} is quenched away
faster than F$_2$ and \Fm{} absorb). That points at the F$_2$-side
absorbers, F$_2$'s compilation value and \Fm's 1971 measurement, and at
KrF$^*$ + F$_2$ quenching, none of them moved here. A proxy for laser
output must saturate: the small-signal net gain quenches the whole
inversion and peaks a factor two low in pressure; the paper's balance is
the honest proxy, and even it falls short of the measurement.

\section{The energy ladder}

v0\_5's ladder (Figure~\ref{fig:ladder}) is v0\_4's with K76 and K76r on
the measurement: the 5p rung \lad{E.PP}\unit{eV} above Kr and
\lad{gap5p}\unit{eV} above \Krs{}, KPS and K27X as quenches, PRAD's
\lad{hv.PRAD}-eV photon, the two stepwise ionizations and the harpoon from
5p; \Krtp{} + \el{} lands on it, leaving \lad{heat.R11p}\unit{eV} to the
gas. KrF(B) lies \lad{DeB}\unit{eV} below the ion pair \Krp{} + \Fm{} and
\lad{BoverKrF}\unit{eV} above Kr + F, and the \lad{wl.G4}-nm photon
carries \lad{hv.G4}\unit{eV} of that.

\begin{figure}[p]
\centering
\ladpicture
\caption{The energy ladder of LIGHT v0\_5, version \lad{version}, in eV
above ground-state Kr and \Ft{} at rest. A bar is a species at its
formation energy and a dashed bar a pair, the sum of two. An arrow is a
row, labelled with its id. It is solid in the default network and dotted
for a band whose default is 0, \lad{zero}. The rows of the optional group
\lad{groups} are dashed. A photon's arrow drops by its quantum, and the red
squiggle below it is the heat its lower state gives the gas as it falls
apart. A red stub is a quench. The energy a row releases heats the gas,
except a photon's, which leaves, and that of \lad{toe}, which goes to the
electrons. KrF(B) and KrF(C) lie \lad{gapBC}\unit{cm^{-1}} apart and are
drawn as a pair of bars that every B and C row meets. The rows
\lad{notdrawn} move nothing between levels and are not drawn.}
\label{fig:ladder}
\end{figure}

\section{The new terms' checks}
\label{sec:checks}

\begin{itemize}
\item The tables as printed: spot values of Jung et al.'s Tables II to IV
are pinned (\texttt{test\_v0\_5.py}), and the sums at 6\unit{eV} are
$\wg{jung.sum.1s5}$ and $\wg{jung.sum.1s3}\unit{cm^2}$.
\item The Maxwellian table of the lumped cross section against a direct
120\,001-point trapezoid at 1, 2 and 5.8\unit{eV}: $2\times10^{-3}$.
\item The solver's population process on \wg{pop.cells} cells to
\wg{pop.emax}\unit{eV}: the forward rate for a Maxwellian against the
table to \wg{pop.k}, the reverse against the Klein-Rosseland table to
\wg{pop.sup}, detailed balance over the gap to \wg{pop.db}.
\item The switch: \texttt{k76 = "kannari"} reproduces v0\_4's reference
point (Table~\ref{tab:ref}'s last column against v0\_4's overview) to four
digits.
\item A \wg{pop.short.t}-ns kinetic pulse with the measured pair inside
the equation: the field carries \wg{pop.short.ratio.lo} to
\wg{pop.short.ratio.hi} of the input and the balances agree to
\wg{pop.short.gap}.
\item v0\_4's checks, unchanged: the wave's field term (atoms to
\wg{hf.atoms}, the marginal ion part to \wg{hf.ions}), the ion-ion limbs
($\wg{ii.k1}$ and $\wg{ii.k1b}\cm{3}$ at 1 amagat), the walls'
\wg{wl.nu}\unit{s^{-1}} at mid-pulse, \texttt{ib.py} by \wg{ib.checks}
checks.
\end{itemize}

\section{Where v0\_5 is not yet reliable}

\begin{itemize}
\item The 5s $\to$ 5p excitation is measured; the two $J=1$ levels are
not, and their proxies are the band's low end. Jung et al.'s absolute
scale is $\pm$\wg{jung.unc.1s5} to \wg{jung.unc.1s3}\,\%, the band's high
end.
\item The trimer's rate at the laser's gas temperature: the one laser
that bears on it finds LIGHT's quenching \x{\wg{kl.k.ratio}} too strong
at \wg{kl.tg}\unit{K}, through a model-dependent constant. The
temperature law is still the largest term of the propagated uncertainty.
\item The F$_2$-side absorbers and KrF$^*$ + F$_2$: the laser's F$_2$
optimum is not reproduced.
\item KPS one level's total loss, K27 a search snippet, the stepwise
ionizations and the 5p harpoon as LIGHT-side rates, the absorbers as a
diagnostic, the wall closure, no propagation: as in v0\_4.
\item Everything in v0\_3's list: the steady EEDF after the pulse, the
first cell, the thermal electrons' grid bias, \Krts{} + F, the trimer above
its measured range, Hakuta's and Burnham and Searles' conflicts, and the
pure-krypton VUV yield at low pump.
\end{itemize}

\appendix
\section{Reaction rates}

v0\_5's rows are v0\_4's (its overview's Appendix A, and v0\_2's) with
these two changed. Section 8 of \texttt{v0\_5/SOURCES.md} says how each
source was read.

\begin{table}[H]
\centering\scriptsize
\renewcommand{\arraystretch}{1.1}
\begin{tabularx}{\linewidth}{@{}>{\raggedright\arraybackslash}p{1.7cm}>{\raggedright\arraybackslash}p{3.9cm}>{\raggedright\arraybackslash}p{3.5cm}>{\raggedright\arraybackslash}X@{}}
\toprule
row & reaction & value & source\\
\midrule
K76 & \el{} + \Krs{} \ra \el{} + \Krpp{} & the lumped measured cross section,
  $\wg{jung.sig6}\unit{cm^2}$ at 6\unit{eV}; $k=\wg{jung.k.116}$ at 1.16 and
  $\wg{jung.k.58}\cm{3}$ at 5.8\unit{eV} (Maxwellian); band k\_5s5p a scale
  \wg{k.5s5p.lo} to \wg{k.5s5p.hi} & Jung et al.~\cite{jung06} Tables I to IV,
  \cite{jung05} Table I; the $J=1$ levels on their same-core metastable's
  set; Kannari's fit~\cite{kannari85,jacob76} as \texttt{k76 = "kannari"}\\
K76r & \el{} + \Krpp{} \ra \el{} + \Krs{} & the Klein-Rosseland reverse, $g$
  12/36 over the gap & detailed balance\\
PRAD (unchanged) & \Krpp{} \ra \Krs{} + $h\nu$ & $\tau=\wg{tau5p}\unit{ns}$ &
  Dzier\.z\k{e}ga et al.~\cite{dzierzega00}; the TALIF review's
  \wg{tau.talif}\unit{ns} for 5p$'$[3/2]$_2$~\cite{gazeli21rev} as a
  cross-check\\
\bottomrule
\end{tabularx}
\caption{The rows v0\_5 changes. Rates in cm$^3$/s.}
\end{table}

\begin{thebibliography}{99}
\footnotesize
\setlength{\itemsep}{1pt}
\bibitem{eden78} Eden, Waynant, Searles and Burnham, Appl.\ Phys.\ Lett.\ 32, 733 (1978).
\bibitem{gadomski92} Gadomski, Xu, Setser and Nelson, Chem.\ Phys.\ Lett.\ 189, 153 (1992).
\bibitem{yu90} Yu and Setser, J.\ Phys.\ Chem.\ 94, 2934 (1990).
\bibitem{xu93} Xu, Gadomski and Setser, J.\ Chem.\ Phys.\ 99, 2591 (1993).
\bibitem{tahoma89} Spectra Technology, Inc., \emph{Experimental Support for KrF Laser Modeling}, final report to Los Alamos National Laboratory, 19 July 1989, OSTI 10136715 (Tables 3-2 and 3-3).
\bibitem{hakuta90} Hakuta, Fujino, Nakayama and Takuma, Jpn.\ J.\ Appl.\ Phys.\ 29, 1194 (1990).
\bibitem{burnham77} Burnham and Searles, J.\ Chem.\ Phys.\ 67, 5967 (1977).
\bibitem{morozov08} Morozov, Heindl, Kr\"ucken, Ulrich and Wieser, J.\ Appl.\ Phys.\ 103, 103301 (2008).
\bibitem{hakuta87} Hakuta, Komori, Mukai and Takuma, J.\ Appl.\ Phys.\ 61, 2113 (1987).
\bibitem{kannari85} Kannari, Obara and Fujioka, J.\ Appl.\ Phys.\ 57, 4309 (1985), Table I and Sec.\ II.B.
\bibitem{jacob76} Jacob and Mangano, Appl.\ Phys.\ Lett.\ 28, 724 (1976).
\bibitem{jung05} Jung, Stone, Boffard, Anderson and Lin, Phys.\ Rev.\ Lett.\ 94, 163202 (2005).
\bibitem{jung06} Jung, Stone, Boffard, Anderson and Lin, Phys.\ Rev.\ A 73, 022722 (2006), Tables I to IV.
\bibitem{gazeli21rev} Gazeli, Lombardi, Aubert, Duluard, Prasanna and Hassouni, Plasma 4, 145 (2021), Table 3.
\bibitem{gazeli21} Gazeli et al., Phys.\ Plasmas 28, 043301 (2021).
\bibitem{dzierzega00} Dzier\.z\k{e}ga, Volz, Nave and Griesmann, Phys.\ Rev.\ A 62, 022505 (2000), Table 1.
\bibitem{kleikamp86} Kleikamp, Peters and Witteman, IEEE J.\ Quantum Electron.\ 22, 219 (1986), Figs.\ 3 and 4, Table I.
\bibitem{berger64} Berger and Seltzer, \emph{Tables of Energy Losses and Ranges of Electrons and Positrons}, NASA SP-3012 (1964), as Kleikamp et al.\ cite it.
\bibitem{shiu77} Shiu and Biondi, Phys.\ Rev.\ A 16, 1817 (1977).
\bibitem{grossl81} Gr\"ossl et al., J.\ Chem.\ Phys.\ 74, 1728 (1981).
\bibitem{tellinghuisen76} Tellinghuisen, Hays, Hoffman and Tisone, J.\ Chem.\ Phys.\ 65, 4473 (1976).
\bibitem{lightv04} \emph{How accurate LIGHT v0\_4 is}, \texttt{LIGHT/docs/accuracy\_overview\_v0\_4.pdf} (2026-10-09), this repository.
\end{thebibliography}

\end{document}
